% Root proof promoted after two complete independent audits; exact candidate retained.
\subsection{A two-rank-one theta curve with a bad five-kernel}

\begin{proposition}[Not every five-quotient is nodal]
\label{prop:root-f1-bad-five-kernel}
\status{Proved; the stronger every-kernel nodality claim is refuted.}
Over $k=\overline{\mathbf F}_2$, let $C$ be the smooth projective
model of
\[
                 y^2+xy=x^5+x^3+x.
\]
Its genus is two and its two-rank is one. Put $Q=(0,0)$ and
$P=(1,\omega)$, where $\omega^2+\omega+1=0$, and identify
$C\subset B=\operatorname{Jac}(C)$ by $R\mapsto[R-Q]$.
The class $D=[P-Q]$ has exact order five. Every maximal isotropic
subgroup $K\subset B[5](k)$ containing $\langle D\rangle$
has at least three points in a fibre of its theta restriction.
There are exactly six such $K$. In particular these polarized
five-quotients have nonnodal theta image.
\end{proposition}
\begin{proof}
The degree-five hyperelliptic model is smooth: the only possible
affine singularity has $x=0$, hence $y=0$, where the derivative
of the defining equation with respect to $x$ is one. The smooth
degree-five completion has one point at infinity and genus two.
With the basis $dx/x,dx$ of regular differentials, the Cartier
operator has matrix $\operatorname{diag}(1,0)$, by the previously
proved genus-two formula. Thus its stable rank is one.

Set $z=1/x$, $v=y/x^3$. Then
\[
             v^2+z^2v=z^5+z^3+z.
\]
The point $Q$ is now infinity, where $z$ has pole order two and
$v$ has pole order five. Consider
\[
              f=v+\omega z^2+z+1.
\]
The hyperelliptic involution sends $v$ to $v+z^2$. Direct expansion
using $\omega^2+\omega=1$ gives
\[
      f\,\iota(f)=(z+1)^5.
\]
At $P$, $z=1,v=\omega$, so $f(P)=0$ and $\iota(f)(P)=1$.
The $z$-map is etale at $P$, since its defining derivative in
$v$ is $z^2=1$. Therefore $f$ has a zero of order five at $P$.
Its sole pole is $Q$, of order five, because the term $v$ has
strictly larger pole order than the other terms. Consequently
\[
                   \operatorname{div}(f)=5P-5Q.
\]
The Abel--Jacobi embedding is injective and $P\ne Q$, so $D$ is
nonzero and has exact order five.

The point $Q$ is Weierstrass. For $P'=(1,\omega+1)=\iota(P)$,
the canonical-fibre relation gives $[P'-Q]=-D$.
Every quotient with kernel containing $D$ therefore identifies
the three distinct points $Q,P,P'$. The symplectic space $B[5](k)$
has dimension four over $\mathbf F_5$. Every line is isotropic,
and maximal isotropic planes containing $\langle D\rangle$
correspond to the lines of the symplectic plane
$\langle D\rangle^\perp/\langle D\rangle$. Their number is
$|\mathbf P^1(\mathbf F_5)|=6$. Five times the source principal polarization descends through
each such kernel to a principal polarization on the quotient.
The proved odd-kernel lemma makes each theta image integral with
normalization $C$; a fibre with three normalization points cannot
be nodal.
\end{proof}

\begin{remark}[Scope of the failed shortcut]
\label{rem:root-f1-bad-five-kernel-scope}
\status{Proved scope statement; no existence theorem is refuted.}
This disproves the stronger shortcut that every five-quotient of
a two-rank-one theta curve is nodal. It does not disprove the
existence of at least one good quotient, or of a suitable incoming
five-neighbour for any fixed target. It closes no main item and
is not counted as progress-review movement on its own.
The norm-function method is consistent with
\citep[Proposition~1.1]{BoxallGrantLeprevost2001}; the full divisor
verification is given above, without a novelty claim.
The bounded search \path{code/root_f1_five_packets.py} checks all
rational points of twenty-four specified curve/field pairs in
$1.321733534001396$ seconds. The separate retained program
\path{code/root_f1_five_packet_function_certificate.py}
checks the displayed function identity without Cantor arithmetic.
No geometric absence claim is inferred from any bounded search.
\end{remark}
